\documentclass[12pt]{article}

\usepackage{exsheets}

%\usepackage[pdf]{pstricks}
%\usepackage{pstricks-add}

%\usepackage{tikz} 
\usepackage[fleqn]{amsmath}
\usepackage{amssymb}
\usepackage{amstext}  
%\usepackage{amsthm} 
\usepackage{array}
\usepackage{mathtools}
\mathtoolsset{showonlyrefs} 
%\usepackage{graphicx,hyperref}
\usepackage{listings}
%\usepackage{indentfirst}
\usepackage[top=1.5cm, bottom=1.5cm, left=.5in, right=.5in]{geometry}
\usepackage{enumerate}
%\usepackage{multicol}

\newcommand{\ssol}{\vspace{3em}}
\newcommand{\lsol}{\vspace{10em}}
\newcommand{\blnk}{{\underline {\hspace{1.5in}}}}

%\usepackage{mcode}

\setcounter{secnumdepth}{0} % Turns off automatic numbering of sections.

\SetupExSheets{headings=block}



\begin{document}
%%%(change to appropriate class and semester)



Justin Boyer\\
Mapping from polynomials with integer coefficients to integers\\

This write up:
\begin{itemize}
\item Defines a mapping to take a polynomial with integer coefficients into the integers.
\item Shows that this mapping preserves addition and multiplication.  \emph{To do: Make it an ring homomorphism: Does it preserve the identity, i.e $ \Phi(1_F) = 1$}?
\item Shows that a polynomial maps to zero if and only if the polynomial is divisible by $x-b$, where $b$ is the base of integer representation we are interested in using (commonly $b=10$).
\end{itemize}


\subsection*{The Mapping $\Phi$}
Let $f(x) = a_n x^n +a_{n-1} x^{n-1} + \cdots + a_0$, where $a_i \in \mathbb{Z}$, for $i=0, 1, 2, \ldots, n$.\\

Define:
$$
\Phi_b(f(x)) = \int f(x) \delta(b-x)dx = f(b)
$$
where $\delta$ is the dirac delta, and $b$ is the base of the integer representation we are mapping into.

So if we were mapping into the base-10 representation of the integers, denoted $\mathbb{Z}_{10}$, we would have:
$$
\Phi_{10}(f(x))  = \int f(x) \delta(10-x) dx = f(10)
$$
It should be stated that $\Phi$ maps the polynomials with integer coefficients to the integers.
\paragraph*{Example}
$f(x)= 5x^4+4x^3+2x+1$, So $\Phi_{10}(f(x))= f(10) = 5(10)^4 + 4(10)^3 + 2(10) + 1 = 54021$  In other words $5x^4+4x^3+2x+1 \xrightarrow[\Phi_{10}]{} 54021$, which means $5x^4+4x^3+2x+1$ goes to $54021$ under $\Phi_{10}$.

\subsection*{$\Phi$ preserves addition and multiplication}
%In order to show that $\Phi$ is a ring homomorphism it is necessary to show that it is closed under addition and multiplication, \emph{but I also need to show $\Phi(1_F) = 1$}.

Given two polynomials with integer coefficients $f_1, f_2$.
\paragraph*{Closed under addition}
\begin{align*}
\Phi_b(f_1(x)+f_2(x)) &= \int (f_1(x) + f_2(x)) \delta(b-x) dx \\
&= \int f_1(x) \delta(b-x) dx + \int f_2(x) \delta(b-x) dx \\
&= f_1(b) + f_2(b)\\
&= \Phi_b(f_1(x)) + \Phi_b(f_2(x))
\end{align*}

\paragraph{Closed under multiplication}
\begin{align*}
\Phi_b(f_1(x) f_2(x)) &= \int f_1(x) f_2(x) \delta(b-x) dx \\
&= f_1(b) f_2(b) - \int f'_1(x) g_2(x)dx \text{ By integrating by parts }\\
&= f_1(b) f_2(b)\\
&= \Phi_b(f_1(x)) \Phi_b(f_2(x))
\end{align*}

Where $\int f'_1(x) g_2(x) dx = 0$ because $g_2(b) = f_2(b)$\footnote{$f_2(b) \in \mathbb{Z}$} when $x=b$, but $g_2(x)$ is zero everywhere else, so the integral has measure zero.


\subsection*{A polynomial maps to zero if and only if the polynomial is divisible by $x-b$}
Now I will show that a polynomial maps to zero if and only if the polynomial is divisible by $x-b$, where $b$ is the base used for the integers.  
\paragraph*{If $x-b$ divides the polynomial then $\Phi \rightarrow 0$}
Take $f(x)$  a polynomial with integer coefficients.  Since $x-b$ divides $f(x)$ we may write $f(x) = (x-b)g(x)$, where $g(x)$ is a nonzero polynomial with integer coefficients.
\begin{align*}
\Phi_b(f(x)) &= \Phi_b(x-b) \Phi_b(g(x)) \\
&= (b-b) g(b) \\
&= 0 
\end{align*}


\paragraph*{If a polynomial maps to zero under $\Phi_b$ then the polynomial is divisible by $x-b$}
Take a non-zero polynomial with integer coefficients, $f(x)$. By the remainder theorem, $f(x)=q(x)(x-c) + f(c)$, for some polynomial with integer coefficients $q(x)$ and $c, f(c) \in \mathbb{Z}$.\\

 Then by problem statement, $\Phi_b(f(x)) = f(b) = 0 $.  So $f(b) = q(b)(b-c) + f(c) = 0$. Because $q(b) \neq 0$ it is possible to write $b-c = \frac{-f(c)}{q(b)}$.\footnote{$q(b) \neq 0$, because if it did then $q(b) = 0 \Rightarrow f(c) = 0, \forall c$ then $q(b) = 0$ to begin with, which implies $f(x) = 0$ to initially.}  A solution exists if $f(c)=0$, well $f(b)=0$ so take $c=b$, which implies $b-b = \frac{f(b)}{q(b)} \Rightarrow 0 = \frac{0}{q(b)}$.  Therefore $c=b$ and if a polynomial maps to 0 under $\Phi_b$ it is divisible by $x-b$


\end{document}


\begin{pspicture}
\psgrid[unit=.05\linewidth, subgriddiv=1, linewidth=.01 gridlabels=0pt, gridcolor=gray, griddots=10](0,0)(-10,-10)(10,10) 
\psaxes[unit=.05\linewidth, linewidth=.1 labelFontSize=\scriptscriptstyle]{<->}(0,0)(-10.5,-10.5)(10.5,10.5)
\end{pspicture}


\begin{tikzpicture}[scale=.5]	
\draw[help lines] (0,0) grid (20,20);
\end{tikzpicture}


